% Adapted from Alain Matthes, including Herbert Voss's completing-square example.
% CC BY-SA 4.0. See ATTRIBUTION.md for provenance and the corrected annotation.
\documentclass[tikz,border=8pt]{standalone}
\usepackage{minimalist-profile}
\usetikzlibrary{calc}
\begin{document}
\begin{tikzpicture}[ms base]
  \node[ms title,text=msColor5,anchor=west] at (0,0.9) {Solve a quadratic equation};
  \node (a1) at (1.7,0) {$3\left(x^2-\frac23\right)=4$};
  \node (a2) at (1.7,-1.15) {$3x^2-2=4$};
  \node (a3) at (1.7,-2.3) {$3x^2=6$};
  \node (a4) at (1.7,-3.45) {$x^2=2$};
  \node (a5) at (1.7,-4.6) {$\sqrt{x^2}=\sqrt2$};
  \node (a6) at (1.7,-5.75) {$|x|=\sqrt2$};
  \node[text=msColor5] (a7) at (1.7,-6.9) {$x=\pm\sqrt2$};
  \foreach \from/\to/\label in {1/2/{expand},2/3/{$+2$},3/4/{$\div3$},4/5/{$\sqrt{\ldots}$},5/6/{$\sqrt{x^2}=|x|$},6/7/{so that}} {
    \draw[ms arrow,draw=msColor5] (a\from.south) -- (a\to.north)
      node[midway,right=\msSpaceLabel,text=msColor5] {\label};
  }
  \node[ms small,text=msMuted,align=left,anchor=west] at (4,-3.45)
    {Isolate the term\\with the variable.};

  \node[ms title,text=msColor1,anchor=west] at (7,0.9) {Complete the square};
  \node (b1) at (10.5,0) {$y=2x^2-3x+5$};
  \node (b2) at (10.5,-2.35)
    {$y=2\left(\textcolor{msColor1}{\underbrace{x^2-\frac32x+\left(\frac34\right)^2}}
      \underbrace{-\left(\frac34\right)^2+\frac52}\right)$};
  \node (b3) at (10.5,-4.6)
    {$y=2\left(\textcolor{msColor1}{\left(x-\frac34\right)^2}+\frac{31}{16}\right)$};
  \node (b4) at (10.5,-6.9)
    {$y=2\textcolor{msColor1}{\left(x-\frac34\right)^2}+\frac{31}{8}$};
  \draw[ms arrow,draw=msColor1] (b1.south) -- (b2.north)
    node[midway,right=\msSpaceLabel,align=left,text=msColor1]
    {Form a binomial square:\\add and subtract $\frac9{16}$ inside.};
  \draw[ms arrow,draw=msColor1] (b2.south) -- (b3.north)
    node[midway,right=\msSpaceLabel,align=left,text=msColor1]
    {$(a-b)^2=a^2-2ab+b^2$};
  \draw[ms arrow,draw=msColor1] (b3.south) -- (b4.north)
    node[midway,right=\msSpaceLabel,text=msColor1] {simplify};
  \node[ms small,text=msMuted,anchor=north west] at (0,-7.5)
    {Real-valued $x$; arrows indicate equivalent algebraic steps.};
\end{tikzpicture}
\end{document}
